% rn-state-and-data.tex — React Native state taxonomy and the right tool for each:
% Context, Redux Toolkit, Zustand/Jotai, TanStack Query, forms, persistence, offline.
% Sources (checked 2026-10-06):
%   tanstack.com/query docs (GitHub TanStack/query): guides/important-defaults (stale by default,
%     gcTime 5 min, retry 3 with exponential backoff) · react-native (onlineManager + NetInfo,
%     focusManager + AppState)
%   github.com/pmndrs/zustand README (no provider; useShallow; persist middleware)
%   github.com/mrousavy/react-native-mmkv README (JSI, fully synchronous, encryption)
%   react.dev/reference/react/useContext (every consumer re-renders on a new value)
%   Not in those sources, marked \unverified: AsyncStorage's Android size cap.
% @labels: area=react-native kind=architecture platform=cross-platform topic=data,architecture discussion=214
% @tags: tanstack-query, redux-toolkit, zustand, jotai, react-context, staletime, optimistic-update, mmkv, asyncstorage, react-hook-form, offline-first, normalisation
\documentclass[8pt]{extarticle}
\usepackage{printup-sheet}

\lstdefinelanguage{TSSheet}{
  morekeywords={const,let,function,return,import,from,export,default,type,
    interface,if,else,new,true,false,null,undefined,async,await,
    useQuery,useMutation,useQueryClient},
  sensitive=true, morecomment=[l]{//}, morecomment=[s]{/*}{*/},
  morestring=[b]", morestring=[b]', morestring=[b]`}

\tikzset{
  q/.style={box, draw=sheetOrange, fill=sheetOrange!8, font=\scriptsize, inner sep=1.5pt,
            minimum height=8mm, text width=21mm},
  tool/.style={box, draw=sheetGreen!70!black, fill=sheetGreen!12, font=\scriptsize,
            inner sep=1.5pt, text width=22mm, minimum height=10mm},
  per/.style={box, font=\scriptsize, inner sep=1.5pt, text width=31mm, minimum height=9mm},
  lbl/.style={font=\tiny, text=black!75, inner sep=1pt, align=center},
}

\begin{document}

\sheettitle{RN state \& data — the taxonomy, and the tool for each kind}{react-native · folio}

\oneliner{``Which state library?'' is the wrong question. Classify first — \textbf{server},
\textbf{navigation}, \textbf{local UI}, \textbf{shared client}, \textbf{form}, \textbf{persisted}
(\emph{secret}?) — each has a best-fit tool; most bugs are one kind kept in another's tool.}

\vspace{3pt}
\noindent\begin{tikzpicture}[sheet, node distance=3mm]
  \node[font=\scriptsize\bfseries, anchor=west] at (-0.3,3.55) {Decide per piece of state — first ``yes'' wins:};
  \node[q] (q1) at (1.0,2.75) {Fetched from a server?};
  \node[q, right=5mm of q1] (q2) {Identifies the screen? (id, tab, filter)};
  \node[q, right=5mm of q2] (q3) {Form input + validation?};
  \node[q, right=5mm of q3] (q4) {Used by one component / subtree?};
  \node[q, right=5mm of q4] (q5) {App-wide, changes rarely? (theme, locale, user)};
  \node[tool, below=4mm of q1] (t1) {\textbf{TanStack Query} / RTK Query — cache, not state};
  \node[tool, below=4mm of q2] (t2) {\textbf{route params} (React Navigation / Expo Router) — ids, not objects};
  \node[tool, below=4mm of q3] (t3) {\textbf{react-hook-form} + schema (zod)};
  \node[tool, below=4mm of q4] (t4) {\textbf{useState / useReducer}; lift up, or compose};
  \node[tool, below=4mm of q5] (t5) {\textbf{Context} — split it, memo the value};
  \node[tool, right=5mm of q5, yshift=-5mm, text width=24mm, minimum height=20mm,
        draw=sheetBlue, fill=sheetBlue!8] (t6) {\textbf{client store}\\[1pt]
        \textbf{Zustand} / \textbf{Jotai}: selectors, tiny API\\[1pt]
        \textbf{Redux Toolkit}: big team, devtools, middleware, strict flow};
  \foreach \a/\b in {q1/q2,q2/q3,q3/q4,q4/q5} \draw[flow] (\a) -- node[lbl, above]{no} (\b);
  \draw[flow] (q5.east) -- node[lbl, above]{no} (q5.east -| t6.west);
  \foreach \a/\b in {q1/t1,q2/t2,q3/t3,q4/t4,q5/t5} \draw[hot] (\a) -- node[lbl, right]{yes} (\b);
  % persistence strip
  \node[font=\scriptsize\bfseries, anchor=west] at (-0.3,0.65) {Must it survive a restart?};
  \node[per, draw=sheetRed, fill=sheetRed!7] at (1.55,-0.05) {\textbf{secret} (token, key)\\Keychain / Keystore:\\\texttt{expo-secure-store}, \texttt{react-native-keychain}};
  \node[per] at (5.75,-0.05) {\textbf{small, hot, read at launch}\\\textbf{MMKV} — sync via JSI, mmap'd, fast; optional encryption};
  \node[per] at (9.95,-0.05) {\textbf{large / queryable / relational}\\SQLite (\texttt{expo-sqlite}, \texttt{op-sqlite}), WatermelonDB};
  \node[per, draw=sheetGrey, fill=black!4] at (14.15,-0.05) {\textbf{simple, legacy}\\AsyncStorage — async, strings, \textbf{unencrypted}};
\end{tikzpicture}

\begin{multicols}{2}
\raggedright

\section{Server state — TanStack Query (v5)}
\begin{itemize}
  \item Cache keyed by \texttt{queryKey} (\texttt{['todo', id]}); dedupes requests;
        \textbf{stale-while-revalidate}: show cached, refetch behind.
  \item \texttt{staleTime} (default \textbf{0}) — how long data is fresh (no refetch on
        mount/focus). \texttt{gcTime} (default \textbf{5\,min}; v4 \texttt{cacheTime}) —
        how long an \emph{unused} entry stays in memory.
  \item \texttt{invalidateQueries} after a mutation: mark stale, refetch active.
        \texttt{retry}: 3, exponential backoff.
  \item RN wiring: \texttt{focusManager} ← \texttt{AppState}, \texttt{onlineManager} ←
        NetInfo. Offline, queries pause; persist the cache with
        \texttt{PersistQueryClientProvider}.
  \item RTK Query: the same inside Redux — \texttt{createApi}, \texttt{providesTags} /
        \texttt{invalidatesTags}.
\end{itemize}

\section{Example — optimistic update with rollback}
\begin{lstlisting}[language=TSSheet]
const qc = useQueryClient();
const save = useMutation({
  mutationFn: (t: Todo) => api.update(t),
  onMutate: async (t) => {
    await qc.cancelQueries({ queryKey: ['todos'] }); // no race
    const prev = qc.getQueryData<Todo[]>(['todos']);
    qc.setQueryData<Todo[]>(['todos'], (old = []) =>
      old.map((x) => (x.id === t.id ? t : x)));
    return { prev };                                  // -> ctx
  },
  onError: (_e, _t, ctx) => qc.setQueryData(['todos'], ctx?.prev),
  onSettled: () => qc.invalidateQueries({ queryKey: ['todos'] }),
});
\end{lstlisting}

\section{Context vs stores — re-render mechanics}
\begin{itemize}
  \item \textbf{Context} re-renders \emph{every} consumer when the value identity
        changes — no selectors. Low-frequency values only; split state from dispatch,
        \texttt{useMemo} the value.
  \item \textbf{Redux Toolkit}: \texttt{createSlice} (reducers ``mutate'' a draft,
        \textbf{Immer} makes the immutable copy), \texttt{configureStore},
        \texttt{createAsyncThunk}. \texttt{useSelector} re-renders when its result
        changes by reference; derive with memoised \texttt{createSelector}.
  \item \textbf{Zustand}: \texttt{create(fn)}, no provider; \texttt{useStore(selector)}
        subscribes to one slice; object selectors need \texttt{useShallow};
        \texttt{persist} middleware.
  \item \textbf{Jotai}: atoms, bottom-up; derived atoms re-run only for their readers.
\end{itemize}

\section{Forms}
\textbf{react-hook-form}: values live outside React state — a keystroke does not re-render
the form; RN inputs go through \texttt{Controller}; validate with a zod resolver.

\columnbreak

\section{Persistence — beyond the picture}
\begin{itemize}
  \item AsyncStorage on Android is SQLite with a default size cap\unverified; MMKV's sync reads
        let settings load at launch without a spinner — its encryption key belongs in
        the Keychain. Secure stores hold \emph{small} values (tokens), not data.
  \item Persist client state selectively and \textbf{version + migrate} its shape.
\end{itemize}

\section{Offline-first \& normalisation}
\begin{itemize}
  \item Local DB is the source of truth; the UI reads it; sync is background.
  \item Writes go to an \textbf{outbox}, replayed on reconnect with
        \textbf{idempotency keys} (a retried POST must not double-charge).
  \item Conflict rule per entity: last-write-wins, server-wins, or field merge.
  \item \textbf{Normalise} shared entities (\texttt{\{ids, entities\}}, RTK
        \texttt{createEntityAdapter}): one update reaches every screen. Query caches
        are per key, \emph{not} normalised — \texttt{setQueryData} or invalidate.
\end{itemize}

\section{Traps}
\begin{itemize}
  \trap{Copying fetched data into Redux/\texttt{useState}: two sources of truth, stale
        forever. Server state belongs in the query cache.}
  \trap{Tokens in AsyncStorage (plain file/SQLite) — use Keychain/Keystore.}
  \trap{\texttt{staleTime: 0} + many mounts = refetch storms; \texttt{gcTime} is
        \emph{not} freshness.}
  \trap{\texttt{useStore()} without a selector re-renders on every change.}
  \trap{Whole objects as route params — stale copies, broken deep links.}
\end{itemize}

\section{Remember}
\textbf{Server $\to$ query cache · screen $\to$ route · local $\to$ useState · rare
global $\to$ Context · hot shared $\to$ store · secret $\to$ Keychain.}


\end{multicols}

\noindent{\footnotesize\color{sheetGrey}\textit{Related:} rn-performance · React Navigation
/ Expo Router · Keychain Services · Core Data / SwiftData (offline) · Combine / Observation}

\end{document}
